% ----------------------------------------------------------------------
% Appendix D: validation protocol and convergence data.
% Fragment to be \input into main.tex after \appendix.
% All numbers from paper/experiments/convergence_results.json and
% paper/draft.md; the stagger ablation reproduced from the library
% (lib/absorbing.py) under the same gold-standard protocol.
% ----------------------------------------------------------------------

\appendix{Validation protocol and convergence data}
\label{ch3:appD:validation}

This appendix records the classical validation protocol behind every
reflection number quoted in Section~\ref{ch3:sec:2d} and the raw convergence
data behind the fitted orders of Section~\ref{ch3:sec:circuit}
[\ref{ch3:fig:convergence}(a--d)].

\subsection{Gold-standard reflection protocol}
\label{ch3:appD:protocol}

Measuring the residual reflection of an absorbing layer against the
absorbing dynamics itself is circular; the gold standard is the
open-domain solution.  The truncated domain has $N=128$ points (seven
grid qubits) with the CPML occupying $n_{\mathrm{pml}}$ points at each
end.  The initial data are a Gaussian velocity bump centred at $j=64$
with width $3$ cells and $w(0)=0$, normalised to
$\lVert z_0\rVert_2=1$.  The reference embeds the \emph{same} truncated
data centred in a $4\times$ hard-wall domain ($N_{\mathrm{ref}}=512$,
embedding offset $\Delta=(N_{\mathrm{ref}}-N)/2=192$) and evolves it
with the closed-domain unitary $e^{-iH_{\mathrm{ref}}T}$.  Both
evolutions are computed by dense matrix exponentials, so the comparison
contains no time-integration error.  The reported error is the
pointwise $\ell_2$ difference on the interior window,
\begin{equation}
  E(T)^{2} \;=\; \sum_{f\in\{v,\,w\}}\;
  \sum_{j=n_{\mathrm{pml}}}^{N-n_{\mathrm{pml}}-1}
  \bigl|\, f_j(T) - f^{\mathrm{ref}}_{j+\Delta}(T) \,\bigr|^{2}.
  \label{ch3:appD:eq:gold}
\end{equation}
Two checks make the protocol sound.  (i)~\emph{Window alignment}: at
$T=0$ the windowed difference in Eq.~\eqref{ch3:appD:eq:gold} evaluates to
exactly zero in floating point, pinning the embedding offset and the
index conventions.  (ii)~\emph{Causal validity}: with $c=1$, $h=1$ the
pulse meets the layer interface at depth $j=n_{\mathrm{pml}}$ and its
reflection is back at the centre by
$T\approx2(64-n_{\mathrm{pml}})\le120$ for $n_{\mathrm{pml}}\ge4$,
whereas the far walls of the reference domain are first felt inside the
comparison window only at $T\approx380$.  The plateau times
$T\in\{100,120,140\}$ therefore capture the complete round-trip
reflection of the truncated domain while the reference is still
uncontaminated.  The reported reflection is the maximum of $E(T)$ over
the three plateau times, and the relative plateau spread
[$(\max-\min)/\mathrm{mean}$] is the flatness diagnostic.  A reduced
instance of the same protocol ($N=64$ versus $256$, $T=20$ and $60$,
asserted $<10^{-3}$) is pinned in the unit-test suite
(Section~\ref{ch3:appD:tests}).

\subsection{Reflection versus layer width and design reflection}
\label{ch3:appD:reflection}

\begin{table}
\caption{Gold-standard true reflection (1D CPML, $N=128$ versus the
$N=512$ hard-wall reference, $m=2$): plateau error
$\max_{T\in\{100,120,140\}}E(T)$ and relative plateau spread.  Upper
block: layer-width sweep at $R_0=10^{-3}$; lower block: design-reflection
sweep at $n_{\mathrm{pml}}=12$.}
\label{ch3:appD:tab:reflection}
\begin{tabular}{llll}
\toprule
$n_{\mathrm{pml}}$ & $R_0$ & plateau error & plateau spread \\
\midrule
4  & $10^{-3}$ & $6.7\times10^{-4}$ & $1.0\times10^{-9}$ \\
6  & $10^{-3}$ & $2.6\times10^{-4}$ & $3.3\times10^{-10}$ \\
8  & $10^{-3}$ & $2.6\times10^{-4}$ & $2.2\times10^{-10}$ \\
12 & $10^{-3}$ & $3.4\times10^{-4}$ & $1.7\times10^{-9}$ \\
16 & $10^{-3}$ & $4.2\times10^{-4}$ & $3.3\times10^{-5}$ \\
\midrule
12 & $10^{-2}$ & $4.5\times10^{-3}$ & $1.5\times10^{-10}$ \\
12 & $10^{-3}$ & $3.4\times10^{-4}$ & $1.7\times10^{-9}$ \\
12 & $10^{-4}$ & $4.3\times10^{-5}$ & $1.9\times10^{-7}$ \\
\bottomrule
\end{tabular}
\end{table}

\ref{ch3:appD:tab:reflection} interprets as follows.  The curve is
the sum of two competing contributions.  At the thinnest layer
($n_{\mathrm{pml}}=4$) the discrete transition error dominates: the
measured reflection sits $9.4\times$ above the discrete wall-return
prediction.  For $n_{\mathrm{pml}}\ge8$ the transition error is
negligible and the reflection is set by the wall return, which
\emph{rises toward the design value from below} as the layer widens:
node sampling of the polynomial profile overestimates
$\int\sigma\,dx$ by a factor $1+O(1/n_{\mathrm{pml}})$ (exactly
$\Sigma\sigma h/\!\int\!\sigma = 1.38$, $1.25$, $1.19$, $1.13$,
$1.09$ for $n_{\mathrm{pml}}=4$--$16$), so the effective round-trip
reflection is
$R_{\mathrm{eff}} = R_0^{\,\Sigma\sigma h/\!\int\!\sigma} =
7.1\times10^{-5}$, $1.7\times10^{-4}$, $2.7\times10^{-4}$,
$4.2\times10^{-4}$, $5.2\times10^{-4}$ across the sweep; the measured
values track this prediction within $5$--$20\%$ for
$n_{\mathrm{pml}}\ge8$.  Thin layers therefore \emph{over}-perform
the nominal $R_0$ (at the price of larger transition error), and
widening the layer converges the reflection up to the design value.
We verified this directly in the {\footnotesize SEISMIC\_CPML}
reference implementation~\cite{komatitsch2007unsplit} (the 2D
isotropic elastic code, unmodified kernels, same gold-standard
protocol, $\alpha=0$, $K=1$): its $P$-wave reflection is likewise
non-monotonic and rises from $n_{\mathrm{pml}}=12$ to $16$
($5.8\times10^{-4}\to6.5\times10^{-4}$, against a predicted
$6.5\times10^{-4}\to7.2\times10^{-4}$), converging up toward the
design $R_0$ and matching its own discrete-sum prediction within
${\sim}10\%$ for $n_{\mathrm{pml}}\ge12$.  Two scheme-dependent
differences: the staggered elastic scheme samples the profile at both
full and half points, so its quadrature ratio is
$1+3/(4n_{\mathrm{pml}})$ rather than our
${\sim}1+3/(2n_{\mathrm{pml}})$; and the minimum sits at
$n_{\mathrm{pml}}\approx12$ there because $S$-wave transition
reflection (whose design value $R_0^{c_p/c_s}$ is far smaller) and the
coarser wavelength resolution keep the transition term dominant
further out.  The effect is invisible in normal use of that code, where
the layer width is fixed at ${\sim}10$ cells.

The miscalibration admits a one-line remedy, which we provide as an
option (\texttt{calibration='discrete'}): rescale the staggered profile
\emph{pair} by a single common factor so that the characteristic
damping sum $\Sigma\,(\sigma_v+\sigma_w)h/2$ equals the design
integral exactly.  The realised reflection is then flat at the design
value, $0.81$--$1.08\,R_0$ across the whole sweep
$n_{\mathrm{pml}}=4$--$16$ [red squares in
\ref{ch3:fig:convergence}(a)], i.e.\ the Collino--Tsogka formula
becomes exact on the grid.  The common factor is essential: rescaling
each profile \emph{separately} gives $\sigma_v$ and $\sigma_w$
different damping functions (the node and midpoint quadratures differ),
an in-layer impedance mismatch that reflects at
$O(1/n_{\mathrm{pml}})$---measured at $1.2\times10^{-2}$ to
$6.9\times10^{-2}$, i.e.\ ten to seventy times \emph{worse} than no
calibration at all.  This remedy applies verbatim to classical CPML
codes.  That the level is
$R_0$-controlled rather than an absolute discretisation floor is
established by the lower block: at $n_{\mathrm{pml}}=12$ the measured
reflection tracks the design value over two decades, with ratios
$E/R_0=0.45$, $0.34$, $0.43$---the Collino--Tsogka formula is
predictive within a factor ${\sim}2.3$ in this discretisation, and
lowering $R_0$ lowers the true reflection commensurately, down to
$4.3\times10^{-5}$ at $R_0=10^{-4}$.  This is the floor behaviour
reported for classical CPML implementations by Komatitsch and
Martin~\cite{komatitsch2007unsplit}.  The plateau is flat to
$\le2\times10^{-9}$ (relative) in six of the eight rows; even the two
most weakly damped rows ($n_{\mathrm{pml}}=16$ and $R_0=10^{-4}$) vary
by only $3.3\times10^{-5}$ and $1.9\times10^{-7}$ of their mean across
$T\in\{100,120,140\}$, i.e.\ the round-trip transient has fully decayed and
the quoted numbers are converged in $T$.

\subsection{Recovery order in the \texorpdfstring{$p$}{p} register}
\label{ch3:appD:porder}

\begin{table}
\caption{Schr\"odingerisation recovery error versus $p$-register size
$n_p$ for the kinked $e^{-|p|}$ profile and the $C^{1}$ cubic profile of
Ref.~\protect\cite{jin2024recovery} [1D CPML benchmark, $N=32$,
$n_{\mathrm{pml}}=8$, $\sigma_{\max}=1$, $T=30$, $p\in[-18,18)$,
$\Delta p=36/2^{n_p}$].  Errors are
$\lVert z_{\mathrm{rec}}-e^{AT}z_0\rVert_2$ with
$\lVert z_0\rVert_2=1$, recovered at the library-default slice
$p^{*}=3\Delta p$ (valid at any $p^{*}>0$ since $H_1\preceq0$ here).
The fitted orders are least-squares slopes of $\log_2 E$ versus $n_p$.}
\label{ch3:appD:tab:profiles}
\begin{tabular}{llll}
\toprule
$n_p$ & $\Delta p$ & $e^{-|p|}$ profile & $C^{1}$ cubic profile \\
\midrule
6  & 0.563  & $4.1\times10^{-2}$ & $1.3\times10^{-2}$ \\
7  & 0.281  & $1.6\times10^{-2}$ & $2.6\times10^{-3}$ \\
8  & 0.141  & $3.6\times10^{-3}$ & $1.7\times10^{-4}$ \\
9  & 0.0703 & $1.5\times10^{-3}$ & $1.3\times10^{-5}$ \\
10 & 0.0352 & $1.1\times10^{-3}$ & $5.9\times10^{-6}$ \\
11 & 0.0176 & $6.3\times10^{-4}$ & $2.4\times10^{-6}$ \\
\midrule
\multicolumn{2}{l}{order in $\Delta p$ (all points)} & 1.23 & 2.64 \\
\multicolumn{2}{l}{order in $\Delta p$ ($n_p=6$--$9$)} & 1.64 & 3.40 \\
\bottomrule
\end{tabular}
\end{table}

\ref{ch3:appD:tab:profiles} gives the recovery error of the full
Schr\"odingerised pipeline (exact per-mode evolution, no Trotter) as a
function of the $p$-register size, for the two warping profiles of
Section~\ref{ch3:sec:schro}.  The errors are absolute per unit initial norm;
by $T=30$ the layer has absorbed all but
$\lVert e^{AT}z_0\rVert_2=1.4\times10^{-3}$ of the state, so resolving
the residual field demands the higher-order profile.  Both profiles
meet or exceed their nominal orders [$O(\Delta p)$ kinked,
$O(\Delta p^{2})$ cubic, Remark~4.9 of Ref.~\cite{jin2024recovery}]:
the full-sweep fits give orders $1.23$ and $2.64$, and the pre-floor
fits ($n_p=6$--$9$) give $1.64$ and $3.40$.  Because the default slice
$p^{*}=3\Delta p$ moves with refinement, we re-evaluated both series at
the fixed slice $p^{*}=1$: the fitted orders become $1.64$ and $3.25$,
so the profile ordering is not a slice artefact.  At $n_p=11$ the cubic
profile is ${\sim}260\times$ more accurate ($2.4\times10^{-6}$ versus
$6.3\times10^{-4}$); the flattening of its last two points marks the
approach to the recovery floor of this configuration.

\subsection{Trotter orders at gate level}
\label{ch3:appD:trotter}

\begin{table}
\caption{Gate-level Trotter error versus step count [1D CPML, $N=8$,
$n_{\mathrm{pml}}=2$, $\sigma_{\max}=1$, $T=4$, $n_p=5$, $p_{\max}=8$;
transpiled circuits on a statevector simulator; errors
$\lVert z_{\mathrm{circ}}-z_{\mathrm{ref}}\rVert_2$ with
$\lVert z_0\rVert_2=1$].  ``vs exact S.''\ uses the exact
Schr\"odingerised evolution (exact per-mode propagation at the same $p$
discretisation) as reference, isolating the splitting error; ``vs
$e^{AT}z_0$'' uses the full matrix-exponential reference and therefore
contains the $p$-discretisation floor, measured independently as
$1.29\times10^{-2}$.  Order~2 symmetrises both Trotter levels (Strang
outer split, palindromic term order inner).}
\label{ch3:appD:tab:trotter}
\begin{tabular}{lllll}
\toprule
 & \multicolumn{2}{c}{order 1} & \multicolumn{2}{c}{order 2} \\
steps & vs exact S. & vs $e^{AT}z_0$ & vs exact S. & vs $e^{AT}z_0$ \\
\midrule
10 & $1.36\times10^{-1}$ & $1.32\times10^{-1}$ & $3.24\times10^{-2}$ & $3.40\times10^{-2}$ \\
20 & $6.55\times10^{-2}$ & $6.27\times10^{-2}$ & $8.02\times10^{-3}$ & $1.46\times10^{-2}$ \\
40 & $3.25\times10^{-2}$ & $3.11\times10^{-2}$ & $2.00\times10^{-3}$ & $1.29\times10^{-2}$ \\
80 & $1.62\times10^{-2}$ & $1.75\times10^{-2}$ & $4.99\times10^{-4}$ & $1.29\times10^{-2}$ \\
\midrule
slope & $-1.02$ & $-0.98$ & $-2.01$ & $-0.44$\textsuperscript{a} \\
\bottomrule
\end{tabular}
\\[2pt]
{\footnotesize \textsuperscript{a}\,Floor-limited: not a meaningful order, see text.}
\end{table}

\ref{ch3:appD:tab:trotter} separates the two error sources of the
circuit pipeline of Section~\ref{ch3:sec:circuit}.  Against the exact
Schr\"odingerised evolution the splitting error is clean: fitted slope
$-1.02$ at order~1 and $-2.01$ at order~2 with both Trotter levels
symmetrised [\ref{ch3:fig:convergence}(d)].  Against the full
$e^{AT}z_0$ reference the order-2 series saturates at
$1.29\times10^{-2}$ from $40$ steps onward---exactly the independently
measured $p$-discretisation floor of this deliberately coarse $n_p=5$
register---so the nominal slope $-0.44$ in that column is
floor-limited rather than a defect of the splitting.  The order-1
series does not saturate within the sweep because its splitting error
stays above the floor.  Enlarging the $p$ register lowers the floor at
the rates of \ref{ch3:appD:tab:profiles}.

\subsection{Stagger ablation}
\label{ch3:appD:stagger}

The damping profiles are sampled at $x_j=(j+s)h$ with the stagger $s$
matched to where each field lives (Section~\ref{ch3:sec:cpml}): for the
forward-difference convention, $v$ at integer nodes ($s=0$) and $w$ at
the half-cells $x_{j-1/2}$ ($s=-1/2$).  Rerunning the gold-standard
protocol ($n_{\mathrm{pml}}=12$, $R_0=10^{-3}$) with the $\sigma_w$
stagger ablated gives \ref{ch3:appD:tab:stagger}: misplacing the
profile by half a cell costs two orders of magnitude in true
reflection.  The stagger is load-bearing, not cosmetic.

\begin{table}
\caption{Stagger ablation under the gold-standard protocol
(Section~\protect\ref{ch3:appD:protocol}; 1D CPML, $n_{\mathrm{pml}}=12$,
$R_0=10^{-3}$): plateau reflection versus the stagger offset $s$ of the
$\sigma_w$ sampling, $x_j=(j+s)h$.}
\label{ch3:appD:tab:stagger}
\begin{tabular}{lll}
\toprule
$\sigma_w$ stagger $s$ & plateau error & degradation \\
\midrule
$-1/2$ (correct half-cell) & $3.4\times10^{-4}$ & $1$ \\
$0$ (unstaggered)          & $4.1\times10^{-2}$ & ${\approx}120$ \\
$+1/2$ (wrong sign)        & $8.1\times10^{-2}$ & ${\approx}235$ \\
\bottomrule
\end{tabular}
\end{table}

\subsection{Test suite}
\label{ch3:appD:tests}

The classical and circuit layers are pinned by $24$ unit tests
(\texttt{tests/test\_absorbing.py}), grouped by the property classes
they assert.  (i)~\emph{Generator structure and spectra}: the
Collino--Tsogka profile amplitude and support;
$\max\operatorname{Re}\operatorname{eig}(A)\le0$ to round-off for the
collapsed 1D, general CFS ($\kappa_{\max}=2$, $\alpha_{\max}=0.05$),
and 2D memory-form generators; the Hermitian split $A=H_1+iH_2$ with
$H_1=-\diag(\sigma)\preceq0$ for the collapsed 1D form and the sponge.
(ii)~\emph{Collapse}: the 1D memory form reproduces the collapsed
diagonal damping to $10^{-10}$.  (iii)~\emph{Absorption and reflection
bounds}: interior energy below $10^{-5}$ where the hard-wall reference
retains more than $0.5$; round-trip amplitude reflection below
$5\times10^{-3}$; 2D CPML and sponge energy decay; and the
gold-standard true reflection below $10^{-3}$ on the reduced instance.
(iv)~\emph{Recovery thresholds}: scalar decay recovered to $10^{-4}$;
the cubic profile beating the kinked profile at equal resolution; 1D
recovery below $2\times10^{-3}$ with a plateau-flatness diagnostic; 2D
sponge recovery below $10^{-3}$ with $\lambda^{+}=0$ asserted; and the
2D CPML threshold of Section~\ref{ch3:sec:lemma} verified in both
directions---accurate above $p^{*}=\lambda^{+}T+1$ and at least ten
times worse below.  (v)~\emph{Circuit-versus-matrix agreement}: the
exact-Hamiltonian-gate path against the numpy per-mode pipeline to
$10^{-8}$ (pinning the QFT/FFT conventions and qubit ordering); the
Trotterised Bell circuit to $5\times10^{-2}$; the separated 2D operator
strings against the block Hamiltonian to $10^{-13}$; the per-step
orders $O(\delta t^{2})$ and $O(\delta t^{3})$ of the projector-aware
Bell circuit and its Strang variant; and the 17-qubit end-to-end
gate-level run to $3\times10^{-2}$.  (vi)~\emph{Guards and
regression}: rejection of out-of-range recovery slices, degenerate
layers, and inconsistent damping vectors; and a regression test pinning
the untouched closed-domain library (the mixed Dirichlet/Neumann
Laplacian stencil).  Every bound above is asserted numerically, so any
change that degrades a spectrum, a reflection, a recovery threshold, or
a circuit convention fails the suite.
